\chapter{Conclusion générale}

Ce projet WebCyber a permis de concevoir, de réaliser et de déployer une
application web sécurisée et conteneurisée, tout en respectant des objectifs
clairs de qualité, de reproductibilité et de sécurité.

\section{Bilan du projet}

Les principaux objectifs ont été atteints :

\begin{itemize}
  \item une application de prise de notes fonctionnelle, avec authentification,
        gestion CRUD, archive, corbeille et recherche;
  \item une architecture conteneurisée composée de Nginx, Flask et PostgreSQL
        orchestrés par Docker Compose;
  \item un déploiement sur AWS avec EC2, Elastic IP, Security Group et Route\,53;
  \item une chaîne HTTPS opérationnelle via Let's Encrypt, avec renouvellement
        automatisé;
  \item une approche de sécurité en profondeur fondée sur six couches : réseau,
        hôte, isolation Docker, proxy/TLS, application et base de données;
  \item une campagne de validation couvrant les tests fonctionnels et l'audit
        de sécurité (Nmap, Nikto, SSL Labs/testssl.sh).
\end{itemize}

Ce rapport montre que la solution développée est cohérente avec le
périmètre académique du PFE tout en conservant des pratiques proches de
l'industrie.

\section{Apports méthodologiques}

Le travail réalisé a apporté plusieurs acquis importants :

\begin{itemize}
  \item maîtrise des mécanismes de conteneurisation et des flux de données
        entre services Docker;
  \item compréhension de la sécurisation d'un déploiement web en production :
        gestion des certificats, configuration de Nginx, règles de pare-feu,
        isolation des services;
  \item capacité à documenter une architecture technique complète, de la
        conception à l'implémentation et au test;
  \item adoption d'une démarche d'audit pragmatique, fondée sur des outils
        standard et sur des conclusions actionnables.
\end{itemize}

\section{Limites et perspectives d'évolution}

Bien que la solution soit robuste pour un usage pédagogique et de démonstration,
plusieurs améliorations restent possibles :

\begin{itemize}
  \item mise en place d'un pipeline CI/CD pour automatiser les tests et le
        déploiement;
  \item supervision et journalisation centralisées (Prometheus, Grafana,
        ELK/CloudWatch);
  \item audit de code statique (bandit, semgrep) et tests de sécurité
        automatisés plus poussés;
  \item renforcement de l'authentification par 2FA ou OAuth;
  \item migration vers une solution d'orchestration plus avancée (Kubernetes,
        ECS) pour améliorer la résilience et l'extensibilité;
  \item sauvegardes automatisées de la base PostgreSQL et procédure de
        restauration.
\end{itemize}

Ces pistes constituent des prolongements naturels pour transformer WebCyber
en une plateforme encore plus professionnelle et adaptée à un contexte de
production.

\section{Conclusion}

WebCyber illustre qu'un projet de fin d'études peut être à la fois simple dans
son périmètre fonctionnel et ambitieux dans son approche architecturale et de
sécurité. La démarche adoptée, centrée sur la reproductibilité et la défense en
profondeur, offre une base solide pour évoluer vers des déploiements plus
complexes et des applications plus critiques.
